\documentclass[11pt]{article}
\usepackage[margin=1in]{geometry}
\usepackage{amsmath,amssymb}
\usepackage{booktabs}
\usepackage{hyperref}

\title{Scenario 1 --- Certainty-Equivalent MPC\\
       for Mobile Charging Station Dispatch\\[2pt]
       \large Formal Model of the Multi-Day Cross-Day Receding Horizon}
\author{Approach 1 (Deterministic Certainty-Equivalent MPC)}
\date{}

\newcommand{\dt}{\Delta t}

\begin{document}
\maketitle

Formal companion to the code in \texttt{Receding\_Horizon/code/}: the cross-day
window MILP is built in \texttt{3\_MCSModel.jl} (\texttt{build\_window\_model}); the
deterministic overnight recharge is \texttt{phase2\_overnight\_charge}; the multi-day
closed loop is in \texttt{4\_MPCLoop.jl} (\texttt{run\_mpc}). Plain-English descriptions
and the data schema are in \texttt{README.md}; a line-by-line code audit is in
\texttt{constraints\_code\_vs\_model.txt}.

The controller is a \textbf{multi-day, cross-day RECEDING horizon}. The loop simulates
$D_{\text{tot}} = n_{\text{days}}+1$ days ($n_{\text{days}}$ reported plus one dropped
\emph{buffer} day). Each day is a \textbf{daytime block} of $n_{\text{day}}$ intervals
($n_{\text{day}}$ inferred from the last work-available interval $+$ a return buffer,
capped at the full day $n_{\text{int}}=96$). Every 15-min step re-solves a MILP over the
\textbf{cross-day window}
\[
  K_{\text{win}} = \bigl[\,k_0,\ \dots,\ \text{end of daytime block }
  \min(D_{\text{tot}},\,\text{day}+\ell)\,\bigr],
\]
with lookahead $\ell = \texttt{lookahead\_days}$, laying the remaining part of today and
the next $\ell$ daytime blocks end to end; only the first interval is applied and the
window recedes. \textbf{Nights are not intervals}: the MCS battery is bridged to full at
each evening boundary inside the MILP, and the realized overnight recharge is scheduled
separately by the deterministic \texttt{phase2\_overnight\_charge}; the CEV battery and
any unfinished work carry over. The buffer day absorbs the CEV energy-neutral wrap-up so
it never lands on a reported day.

The activity powers are a \textbf{fixed} calibrated estimate (Fork B); the MILP plans on
the mean $\mu$, and the stochastic plant samples per-interval powers from
$\mathcal{N}(\mu,\sigma)$ truncated at $0$ (idle $\sigma=0$). The model is not re-fit
online.

\paragraph{Window geometry.} For a global interval index $k$, let
$\mathrm{wd}(k)=((k-1)\bmod n_{\text{day}})+1$ be its position \emph{within} its day and
$\mathrm{dayof}(k)=\lfloor(k-1)/n_{\text{day}}\rfloor+1$ its day. Daily profiles (price,
carbon, work availability) repeat each day and are indexed by $\mathrm{wd}(k)$. The
\emph{evening} set $\mathcal{K}^{\text{eve}}=\{k\in K_{\text{win}}:\mathrm{wd}(k)=
n_{\text{day}}\}$ collects each day's end-of-shift interval; interior evenings
$\mathcal{K}^{\text{eve}}\setminus\{\max K_{\text{win}}\}$ carry the overnight MCS reset.

%======================================================================
\section{Sets and indices}
\begin{tabular}{ll}
$m \in \mathcal{M}$ & mobile charging stations (MCS) \\
$e \in \mathcal{E}$ & construction EVs (excavators) \\
$i,j \in \mathcal{N}$ & nodes; $\mathcal{N}_g$ grid, $\mathcal{N}_c$ sites \\
$a \in \mathcal{B}$ & activities $=\{\text{dig},\text{load},\text{trv},\text{idle}\}$ \\
$k$ & global interval in the cross-day window $K_{\text{win}}$ \\
$\mathrm{wd}(k),\mathrm{dayof}(k)$ & within-day index / day of $k$ \\
$\mathcal{K}^{\text{eve}}$ & evening (end-of-shift) intervals in the window \\
$\mathcal{D}$ & day-blocks present in the window \\
$\mathcal{T}_b$ & interval boundaries of the window \\
\end{tabular}

\medskip
$n_{\text{day}}\le n_{\text{int}}=96$ is the \textbf{daytime block} at $\dt=0.25$\,h,
inferred from work availability (e.g.\ $40$ intervals, 08:00$\to$18:00); the overnight
remainder $[n_{\text{day}}{+}1,\dots,n_{\text{int}}]$ is handled by Phase~2. $A_{i,e}=1$
iff CEV $e$ works at site $i$. Work availability $R_{i,e,k}=R^{\text{day}}_{i,e,\mathrm{wd}(k)}$
is $0$ outside the shift, but the MCS/CEVs may charge and the MCS may travel at any
daytime interval.

%======================================================================
\section{Parameters (selected)}
\begin{tabular}{ll}
$\dt$ & interval length (h), $0.25$ \\
$\lambda^{\text{buy}}_{k},\lambda^{\text{CO2}}_{k}$ & price (\$/kWh) / carbon intensity \\
$\overline{SOE}_\bullet,\underline{SOE}_\bullet,SOE^{\text{ini}}_\bullet$ & energy cap / floor / start (MCS, CEV) \\
$CH_m,DCH_m,DCH^{\text{plug}}_m,C^{\text{plug}}_m,\eta_m,CH^{\text{CEV}}_e$ & MCS / CEV charge-discharge parameters \\
$p_a$ & activity power (kW); fixed calibrated mean, MILP plans on it \\
$H^{\text{dig}}_{i,d},H^{\text{load}}_{i,d}$ & \textbf{per-day} dig / load quota (h) at site $i$, day $d$ \\
$r^{\text{dig}}_i,r^{\text{load}}_i$ & work still remaining at site $i$ carried into the window \\
$\tau_{i,j},k_{\text{trv}},R_{i,e,k}$ & travel time / energy / work availability \\
$\rho_{\text{miss}},\rho_{\text{lab}},\lambda^{\text{NC}},\lambda^{\text{OP}}$ & penalties \\
$\text{scale},\kappa=4,r$ & precedence ratio / travels-per-work ($\kappa_{\text{wt}}$) / rest cap \\
$f_{\text{oc}},\Delta_{\text{dd}},\epsilon$ & overcharge frac / drawdown floor (kWh) / terminal tol \\
\end{tabular}

%======================================================================
\section{Decision variables}
\textbf{Continuous ($\ge0$):} $P^{\text{ch}}_{m,i,k}$, $P^{\text{dch}}_{m,i,k}$,
$P^{\text{MCS-CEV}}_{m,i,e,k}$, $P^{\text{work}}_{i,e,k}$, totals
$P^{\text{ch,tot}}_{m,k},P^{\text{dch,tot}}_{m,k}$, travel energy $L_{m,i,j,k}$, state
$SOE^{\text{MCS}}_{m,t},SOE^{\text{CEV}}_{e,t}$, peaks $P^{\text{NC}},P^{\text{OP}}$, and slacks
$s^{\text{miss,dig}}_{i,d},s^{\text{miss,load}}_{i,d}$ (per site \& day-block) plus the soft-fallback
slacks $s^{\text{prec}}_{i,k},s^{\text{pace}\pm}_{e,k},s^{\text{term}}_{e,d}$ (active only under the
matching \texttt{soft\_*} lever).

\smallskip
\textbf{Binary:} $u_{e,i,a,k}$, $\mu_{i,e,k}$, $\rho_{m,i,e,k}$, $z_{m,i,k}$,
$g^{\text{ch}}_{m,i,k}$, $x_{m,i,j,k}$, $y_{m,i,j,k}$, $\beta^{\text{arr}},\beta^{\text{dep}}$.

%======================================================================
\section{Objective (Eq.\ 1)}
\begin{align}
\min\ &\sum_{m,k}\lambda^{\text{buy}}_{\mathrm{wd}(k)}P^{\text{ch,tot}}_{m,k}\dt
      +\sum_{m,k}\tfrac{c_{\text{ton}}}{1000}\lambda^{\text{CO2}}_{\mathrm{wd}(k)}P^{\text{ch,tot}}_{m,k}\dt
      +\rho_{\text{miss}}\!\!\sum_{i\in\mathcal{N}_c,d\in\mathcal{D}}\!\!\bigl(s^{\text{miss,dig}}_{i,d}+s^{\text{miss,load}}_{i,d}\bigr)\notag\\
     &+\lambda^{\text{NC}}P^{\text{NC}}+\lambda^{\text{OP}}P^{\text{OP}}
      +\rho_{\text{lab}}\dt\!\!\sum_{m,i,j,k}\!\! y_{m,i,j,k}
      + W_{\text{p}}\!\sum_{i,k}\! s^{\text{prec}}_{i,k}
      + W_{\text{c}}\!\sum_{e,k}\!(s^{\text{pace}+}_{e,k}{+}s^{\text{pace}-}_{e,k})
      + W_{\text{t}}\!\sum_{e,d}\! s^{\text{term}}_{e,d}.
\end{align}
Price/carbon use the within-day index $\mathrm{wd}(k)$. The soft-fallback penalties
$(W_{\text{p}},W_{\text{c}},W_{\text{t}})=(800,100,150)$ are inert unless the matching lever is
set (the slacks are otherwise pinned to $0$). The overnight recharge cost is reported
separately by Phase~2 and is \emph{not} in this objective.

%======================================================================
\section{Constraints}

\paragraph{Power aggregation \& flow (HARD).}
$P^{\text{ch,tot}}_{m,k}=\sum_{i\in\mathcal{N}_g}P^{\text{ch}}_{m,i,k}$,
$P^{\text{dch,tot}}_{m,k}=\sum_{i\in\mathcal{N}_c}P^{\text{dch}}_{m,i,k}$;
$P^{\text{dch}}_{m,i,k}=\sum_e P^{\text{MCS-CEV}}_{m,i,e,k}\le DCH_m z_{m,i,k}$;
$P^{\text{MCS-CEV}}_{m,i,e,k}\le DCH^{\text{plug}}_m\rho_{m,i,e,k}$;
$\sum_m P^{\text{MCS-CEV}}_{m,i,e,k}\le CH^{\text{CEV}}_e\mu_{i,e,k}$.
Grid exclusivity: $P^{\text{ch}}_{m,i,k}\le CH_m g^{\text{ch}}_{m,i,k}$,
$g^{\text{ch}}\le z$, $\sum_m g^{\text{ch}}_{m,i,k}\le 1$.

\paragraph{State of energy (HARD).}
\begin{align}
SOE^{\text{MCS}}_{m,k+1}&=SOE^{\text{MCS}}_{m,k}+\eta_m P^{\text{ch,tot}}_{m,k}\dt-\tfrac{1}{\eta_m}P^{\text{dch,tot}}_{m,k}\dt-L^{\text{tot}}_{m,k}, && k\notin\mathcal{K}^{\text{eve,int}},\\
SOE^{\text{MCS}}_{m,k+1}&=SOE^{\text{MCS,ini}}_m, && k\in\mathcal{K}^{\text{eve,int}}\ \text{(overnight reset)},\\
SOE^{\text{CEV}}_{e,k+1}&=SOE^{\text{CEV}}_{e,k}+\textstyle\sum_{m,i}P^{\text{MCS-CEV}}_{m,i,e,k}\dt-\sum_i P^{\text{work}}_{i,e,k}\dt, && \text{(carries across nights)}\\
\underline{SOE}_\bullet&\le SOE_\bullet\le\overline{SOE}_\bullet .
\end{align}
Here $\mathcal{K}^{\text{eve,int}}=\mathcal{K}^{\text{eve}}\setminus\{\max K_{\text{win}}\}$ are the
interior evenings. The MCS flow recursion is \emph{skipped} across each night and replaced by
the reset bridge (the realized recharge is Phase~2); the CEV battery carries over unchanged.

\paragraph{Daily CEV neutrality \& health (Eq.\ 8b, HARD; soft under \texttt{soft\_term}).}
At \emph{every} evening boundary $k\in\mathcal{K}^{\text{eve}}$ each CEV must return to (at least)
its start-of-day level, so each reported day is energy-neutral:
\begin{equation}
SOE^{\text{CEV}}_{e,k+1}\ \ge\ SOE^{\text{CEV,ini}}_e-\epsilon,\qquad k\in\mathcal{K}^{\text{eve}}.
\end{equation}
A \emph{floor} (not Avik's exact equality) because a CEV cannot discharge: under the stochastic
plant an equality can become unrecoverable once a CEV drifts high. Under \texttt{soft\_term} this
becomes a penalised two-sided slack $|SOE^{\text{CEV}}_{e,k+1}-SOE^{\text{CEV,ini}}_e|\le
s^{\text{term}}_{e,\mathrm{dayof}(k)}$. The MCS has \emph{no} in-MILP terminal equality --- the
overnight bridge above plus Phase~2 close its cycle.

\paragraph{Keep-up reserve, anti-hoarding cap \& health floor (per day).}
Three state bounds keep the single MCS shuttling and each daily terminal recursively reachable:
\begin{itemize}
\item \textbf{Keep-up reserve (HARD, relaxable):} for each day a backward-built lower bound on
  $SOE^{\text{CEV}}_{e,t}$ from that day's evening deadline (accounting for the intervals the MCS
  must leave early to park at the grid, a share-derated recovery rate, and an idle-drain factor),
  so the terminal is always still reachable.
\item \textbf{Anti-hoarding cap (HARD, always safe):}
  $SOE^{\text{CEV}}_{e,t}\le SOE^{\text{CEV,ini}}_e+f_{\text{oc}}(\overline{SOE}^{\text{CEV}}_e-
  SOE^{\text{CEV,ini}}_e)$ for $t\ne$ start --- no charging far ahead, so the MCS cannot camp at
  the near CEV. An upper bound below $\overline{SOE}$, so it never causes infeasibility.
\item \textbf{Daytime health floor (HARD, relaxable):}
  $SOE^{\text{CEV}}_{e,t}\ge SOE^{\text{CEV,ini}}_e-\Delta_{\text{dd}}$ at daytime boundaries ---
  forbids any CEV draining too far, forcing the MCS to keep both healthy.
\end{itemize}
The reserve and health floor are dropped when the loop re-solves with \texttt{soft\_reserve}
after a hard infeasibility.

\paragraph{Plugging, presence \& routing (Eq.\ 10, HARD).}
$\rho\le A$, $\rho\le z$, $\sum_e \rho_{m,i,e,k}\le C^{\text{plug}}_m$. Presence partition
$\sum_i z_{m,i,k}+\sum_{i\ne j}y_{m,i,j,k}=1$; $\beta^{\text{dep}}_{m,i,k}=\sum_j x_{m,i,j,k}$;
$\beta^{\text{arr}}$ from finishing trips (or a carried-in arrival);
$\beta^{\text{arr}}-\beta^{\text{dep}}=z_k-z_{k-1}$; $\beta^{\text{arr}}+\beta^{\text{dep}}\le1$;
flow balance (generalised for the carried-in start). Trip on $i\!\to\!j$ lasts $\tau_{i,j}$ and
costs $L=k_{\text{trv}}\dt\,y$. \textbf{(10e)} MCS parked at a grid node at \emph{every} evening
(each day's end-of-shift): $\sum_{i\in\mathcal{N}_g}z_{m,i,k}=1,\ k\in\mathcal{K}^{\text{eve}}$.

\paragraph{Activity scheduling (Eq.\ 11, HARD).}
$\sum_a u_{e,i,a,k}=A_{i,e}$; $u\le A$; $\mu_{i,e,k}\le u_{e,i,\text{idle},k}$;
$P^{\text{work}}_{i,e,k}\le R_{i,e,k}A_{i,e}(1-\mu_{i,e,k})$;
$P^{\text{work}}_{i,e,k}=\sum_a p_a u_{e,i,a,k}$ (Eq.\ 5e), with $p_{\text{idle}}=0$.
(The code uses this 4-activity encoding; Avik's equivalent 3-activity form is
$\sum_a u\le1$, $\sum_a u+\mu\le1$.)

\paragraph{Work quota (Eq.\ 12c) --- per-day, cumulative, two-sided.}
Work is a \textbf{per-day} schedule. For each day-block $d\in\mathcal{D}$ let
$K_{\le d}=\{k\in K_{\text{win}}:\mathrm{dayof}(k)\le d\}$ and let the cumulative target
$T^{\text{dig}}_{i,d}=r^{\text{dig}}_i+\sum_{d'\in\mathcal{D},\,d_0<d'\le d}H^{\text{dig}}_{i,d'}$
(likewise load), where $r^{\text{dig}}_i$ carries in the current day's remaining requirement and
$d_0$ is the window's first day. With $D^{\text{dig}}_{i,d}=\dt\sum_{e,k\in K_{\le d}}
u_{e,i,\text{dig},k}$:
\begin{align}
s^{\text{miss,dig}}_{i,d}&\ge T^{\text{dig}}_{i,d}-D^{\text{dig}}_{i,d}, &&\text{(SOFT lower; shortfall rolls over)}\\
D^{\text{dig}}_{i,d}&\le T^{\text{dig}}_{i,d}, &&\text{(HARD upper; \emph{no working ahead})}
\end{align}
and identically for loading. The hard upper cap stops a day from borrowing a future day's quota
now that the window sees tomorrow; unfinished earlier work can still be caught up (same
cumulative budget).

\paragraph{Precedence (Eq.\ 12d, HARD; soft under \texttt{soft\_prec}).}
Raw interval counts, \emph{per day-block} (counters restart each morning; carried realized work
seeds only the current day): for $k$ in day $d=\mathrm{dayof}(k)$ with day-start $b(k)$,
$\sum_{e,\,b(k)\le\tau\le k}u_{e,i,\text{load},\tau}\le\text{scale}\sum_{e,\,b(k)\le\tau\le k}
u_{e,i,\text{dig},\tau}+s^{\text{prec}}_{i,k}$ (with the day-1 seed added on both sides).

\paragraph{Rest rule (Eq.\ 12e, HARD).}
With $r=\lceil t_{\text{rest}}/\dt\rceil$ (=4), over any $(r{+}1)$-window lying inside one day-block
$\sum_{k'=k}^{k+r}\sum_{a\in\{\text{dig,load,trv}\}}u_{e,i,a,k'}\le r$, \emph{plus a seam} seeded
from the current day's applied Work/Break history so a work-run cannot leak across the every-
15-min re-solves. Counters restart each morning (a night is a long break).

\paragraph{Travel pacing (Eq.\ 13, HARD; soft under \texttt{soft\_pace}).}
With $\kappa=\kappa_{\text{wt}}=4$, \emph{per day-block} for each $(i,e)$, on cumulative travel $V$
vs cumulative useful work $W$ (dig+load, seeded with the current day's applied counts):
\begin{equation}
W(k)-\kappa-s^{\text{pace}-}_{e,k}\ \le\ \kappa\,V(k)\ \le\ W(k)+s^{\text{pace}+}_{e,k}.
\end{equation}

\paragraph{Demand peaks (E1, HARD).}
$P^{\text{NC}}\ge$ carried-in peak and $\ge\sum_m P^{\text{ch,tot}}_{m,k}$ for all $k$;
$P^{\text{OP}}$ likewise on on-peak $k$ only.

%======================================================================
\section{Closed loop, overnight Phase 2 \& graceful fallback}
The MILP above is re-solved every 15 minutes over the receding cross-day window; only interval
$k_0$ is applied, then the real state (SOE, MCS position including mid-drive trips, per-site
remaining work, demand peaks, and the current day's applied-activity history) is advanced and the
window slides. The activity powers $p_a$ are the mean of a \emph{fixed} calibrated estimate
$\mathcal{N}(\mu,\sigma)$; the plant draws each interval's realized power from
$\mathcal{N}(\mu,\sigma)$ (truncated at $0$, idle pinned to $0$); the model is not re-fit online.

\paragraph{Overnight Phase 2 (deterministic).} After each daytime block the MCS is refilled from
its end-of-day $SOE^{\text{MCS}}$ back to $SOE^{\text{MCS,ini}}_m$ over the overnight intervals
$[n_{\text{day}}{+}1,\dots,n_{\text{int}}]$, greedily choosing the cheapest overnight prices first
(subject to $\eta_m CH_m$ per interval). This is a closed-form schedule, not part of the MILP.

\paragraph{Graceful fallback.} If a window is hard-infeasible (the daily terminal and keep-up
reserve are jointly unreachable from a drifted state), it is re-solved \emph{once} with the
terminal made soft ($s^{\text{term}}$) and the reserve/health floor dropped
(\texttt{soft\_reserve}), so the MCS still charges toward target rather than freezing; any residual
shortfall surfaces as a penalty. Only if that also fails does the plant \emph{hold} (idle break).

\paragraph{Buffer day.} The loop simulates $D_{\text{tot}}=n_{\text{days}}+1$ days and drops the
last; KPIs and figures use only days $1,\dots,n_{\text{days}}$, so the terminal wrap-up never lands
on a reported day.

\end{document}
